\documentclass[border=6pt]{standalone}
\usepackage{amsmath}
% Estilo común de los esquemáticos (circuitikz)
\usepackage[siunitx, european resistors, american inductors, straightvoltages]{circuitikz}
\usepackage{siunitx}
\sisetup{output-decimal-marker={,}}
\ctikzset{bipoles/length=0.95cm, diodes/scale=0.85, resistors/scale=0.85, capacitors/scale=0.85}
\tikzset{font=\footnotesize, lbl/.style={font=\footnotesize, inner sep=1pt}, net/.style={font=\footnotesize\bfseries, inner sep=1pt}}
% Regulador de tres terminales: \reg{(x,y) del pin IN}{nombre}{pin IN}{pin OUT}{pin ADJ}; caja de 2,6 x 1,4
\newcommand{\reg}[5]{%
  \draw[thick] ($#1+(0,-0.7)$) rectangle ++(2.6,1.4);
  \node[font=\footnotesize\bfseries] at ($#1+(1.3,0.32)$) {#2};
  \node[font=\tiny, anchor=west] at ($#1+(0.05,-0.15)$) {IN};
  \node[font=\tiny, anchor=east] at ($#1+(2.55,-0.15)$) {OUT};
  \node[font=\tiny, anchor=south] at ($#1+(1.3,-0.68)$) {ADJ};
  \node[font=\tiny, anchor=south east] at ($#1+(-0.02,0.02)$) {#3};
  \node[font=\tiny, anchor=south west] at ($#1+(2.62,0.02)$) {#4};
  \node[font=\tiny, anchor=north west] at ($#1+(1.33,-0.72)$) {#5};
}
% etiqueta apilada: nombre arriba, valor abajo
\newcommand{\lab}[2]{\shortstack[l]{#1\\#2}}

\usetikzlibrary{calc}
\begin{document}
\begin{circuitikz}
\node[font=\Large\bfseries, anchor=west] at (-6.4,5.2) {Fuente $\pm$16 V / +5 V para los 4 PID del levitador --- esquemático completo};
\node[font=\small, anchor=west] at (-6.4,4.5) {Las redes con el mismo nombre ($V_{RAW+}$, $V_{RAW-}$, +16 V, GND) están unidas. Brahiam Alberto Campiño Romero, Universidad de Ibagué.};
\begin{scope}[shift={(0,0)}]% ---------- Etapa 1 (v2): red, transformador 18-0-18 V / 3 A, puente rectificador y filtros ----------
\draw (-6.2,1.5) node[ocirc, label=left:{\scriptsize L}]{} to[fuse, l=$F_1$ T1 A] (-3.8,1.5) to[spst, l=$S_1$] (-2.0,1.5) -- (-1.6,1.5);
\draw (-6.2,-1.5) node[ocirc, label=left:{\scriptsize N}]{} -- (-1.6,-1.5);
\node[font=\scriptsize, align=center] at (-5.0,0) {115 V\\60 Hz};
\draw (-1.6,1.5) to[L] (-1.6,-1.5);
\draw (0.0,1.5) to[L] (0.0,0) to[L] (0.0,-1.5);
\draw[thick] (-0.92,1.4) -- (-0.92,-1.4); \draw[thick] (-0.68,1.4) -- (-0.68,-1.4);
\node[font=\footnotesize\bfseries] at (-0.8,2.4) {$T_1$ 18-0-18 V};
\node[font=\scriptsize] at (-0.8,1.95) {115 V, 3 A (108 VA)};
\draw[dashed, gray] (1.0,-3.0) -- (1.0,3.6); \node[gray, font=\scriptsize, anchor=south east] at (0.9,3.0) {fuera de la placa}; \node[gray, font=\scriptsize, anchor=south west] at (1.1,3.0) {placa};
% secundario: extremos A y B con fusible, toma central a GND
\draw (0,1.5) -- (1.4,1.5) node[ocirc]{} node[above, font=\tiny]{J1-1} to[fuse, l={$F_2$ T3,15 A}] (3.6,1.5) node[circ](A){} node[above right, net]{A};
\draw (0,0) -- (1.4,0) node[ocirc]{} node[above, font=\tiny]{J1-2} -- (2.0,0) node[ground]{};
\node[font=\tiny] at (2.6,0.25) {toma central};
\draw (0,-1.5) -- (1.4,-1.5) node[ocirc]{} node[below, font=\tiny]{J1-3} to[fuse, l_={$F_3$ T3,15 A}] (3.6,-1.5) node[circ](B){} node[below right, net]{B};
% puente: D1, D2 hacia V_RAW+ ; D3, D4 desde V_RAW-
\draw (A) -- (4.4,1.5) to[D*, l=$D_1$] (4.4,3.4) -- (6.6,3.4) node[circ](vp){};
\draw (B) -- (5.4,-1.5) to[D*, l_=$D_2$] (5.4,3.4);
\draw (4.4,1.5) node[circ]{} to[D*, invert, l_=$D_3$] (4.4,-3.4) -- (6.6,-3.4) node[circ](vn){};
\draw (5.4,-1.5) node[circ]{} -- (5.4,-1.0) to[D*, invert, l=$D_4$] (5.4,-3.4);
\draw (4.4,-3.4) node[circ]{}; \draw (5.4,-3.4) node[circ]{}; \draw (5.4,3.4) node[circ]{};
\node[font=\scriptsize] at (3.2,-2.6) {$D_1$--$D_4$: 1N5408};
% filtros
\draw (vp) to[eC, l_=$C_1$, a^={\SI{15000}{\micro\farad}/50 V}] (6.6,1.2) node[ground]{};
\draw (vp) -- (9.6,3.4) node[circ]{} to[C, l_=$C_3$, a^=\SI{100}{\nano\farad}] (9.6,1.8) node[ground]{};
\draw (9.6,3.4) -- (13.2,3.4) node[ocirc]{} node[right, net]{$V_{RAW+}$ (TP1)};
\draw (vn) to[eC, invert, l_=$C_{2A}$, a^={\SI{6800}{\micro\farad}/50 V}] (6.6,-5.4) node[ground]{};
\draw (vn) -- (10.2,-3.4) node[circ]{} to[eC, invert, l_=$C_{2B}$, a^={\SI{6800}{\micro\farad}/50 V}] (10.2,-5.4) node[ground]{};
\draw (10.2,-3.4) -- (12.6,-3.4) node[circ]{} to[C, l_=$C_4$, a^=\SI{100}{\nano\farad}] (12.6,-5.0) node[ground]{};
\draw (12.6,-3.4) -- (13.2,-3.4) node[ocirc]{} node[right, net]{$V_{RAW-}$ (TP2)};
\node[font=\scriptsize, anchor=west] at (6.2,-6.3) {$C_{2A}$ y $C_{2B}$ en paralelo = 13\,600~µF (positivo hacia GND)};
\end{scope}
\node[font=\bfseries, anchor=west] at (-6.4,-3.4) {Riel +16 V (limitador $R_{SC1}$--$Q_1$--$Q_2$ y regulador U1)};
\begin{scope}[shift={(-6.0,-5.8)}]% ---------- Etapa 2: limitador de corriente + regulador LM317 de +16 V ----------
% y = 0: riel; y = -4: GND
\draw (0,0) node[ocirc]{} node[left, net]{$V_{RAW+}$} -- (1.0,0) node[circ](a){} to[R, l=$R_{SC1}$, a_={\SI{1.8}{\ohm} 1 W}] (3.4,0) node[circ](ve){} node[above, font=\tiny]{$V_{E1}$} -- (3.9,0);
\draw (3.9,0) node[pnp, rotate=90, anchor=E](Q1){};
\node[lbl, anchor=south] at ($(Q1.B)+(0,1.0)$) {$Q_1$ TIP32C};
\draw (Q1.C) -- (6.0,0) node[circ](vin){} node[above, font=\tiny]{$V_{IN1}$} -- (7.0,0);
\draw (1.6,-1.6) node[pnp, xscale=-1, anchor=B](Q2){};
\draw (Q2.B) -- (2.6,-1.6) -- (2.6,-0.8) -| (ve);
\draw (Q2.E) |- (a);
\node[lbl, anchor=east] at ($(Q2.B)+(-1.0,0)$) {\lab{$Q_2$}{BC557}};
\draw (Q2.C) |- (Q1.B |- 0,-2.5) node[circ](vb){};
\draw (Q1.B) -- (vb) node[right, font=\tiny]{$V_{B1}$};
\draw (vb) to[R, l_={\lab{$R_3$}{\SI{820}{\ohm} 2 W}}] (Q1.B |- 0,-4);
\draw (vin) to[C, l_={\lab{$C_5$}{\SI{100}{\nano\farad}}}] (6.0,-4);
\reg{(7.0,0)}{U1 LM317T}{3}{2}{1}
\draw (9.6,0) -- (10.4,0) node[circ](out){};
\draw (out) -- (10.4,1.4) to[D*, l_=$D_3$ 1N4007] (6.6,1.4) -- (6.6,0) node[circ]{};
\draw (8.3,-0.7) -- (8.3,-1.6) node[circ](adj){} node[above left, font=\tiny]{ADJ1};
\draw (adj) to[R, l_={\lab{$R_7$}{\SI{1.3}{\kilo\ohm} 1 \%}}] (8.3,-2.8) to[vR, l_={\lab{$RV_1$}{\SI{200}{\ohm} 3296}}] (8.3,-4);
\draw (adj) -- (9.6,-1.6) node[circ]{} to[eC, l={\lab{$C_7$}{\SI{10}{\micro\farad}/25 V}}] (9.6,-4);
\draw (out) -- (11.2,0) node[circ]{} to[R, l={\lab{$R_5$}{\SI{120}{\ohm} 1 \%}}] (11.2,-1.6) -- (9.6,-1.6);
\draw (11.2,-1.6) node[circ]{} -- (13.2,-1.6) to[D*, l_={\lab{$D_4$}{1N4148}}] (13.2,0) node[circ]{};
\draw (out) -- (21.0,0) node[ocirc]{} node[right, net]{+16 V};
\draw (15.0,0) node[circ]{} to[eC, l={\lab{$C_9$}{\SI{10}{\micro\farad} tánt.}}] (15.0,-4);
\draw (16.7,0) node[circ]{} to[C, l={\lab{$C_{10}$}{\SI{100}{\nano\farad}}}] (16.7,-4);
\draw (18.3,-4) to[D*, l_={\lab{$D_7$}{1N4007}}] (18.3,0) node[circ]{};
\draw (19.9,0) node[circ]{} to[R, l={\lab{$R_{12}$}{\SI{2.2}{\kilo\ohm}}}] (19.9,-2.2) to[leDo, l={\lab{LED1}{verde}}] (19.9,-4);
\draw (1.0,-4) -- (21.0,-4) node[ocirc]{} node[right, net]{GND};
\foreach \x in {6.0,8.3,9.6,15.0,16.7,18.3,19.9} \draw (\x,-4) node[circ]{};
\draw (Q1.B |- 0,-4) node[circ]{};
\draw (1.0,-4) node[ground]{};
\end{scope}
\node[font=\bfseries, anchor=west] at (-6.4,-10.7) {Riel $-$16 V (limitador $R_{SC2}$--$Q_3$--$Q_4$ y regulador U2)};
\begin{scope}[shift={(-6.0,-13.1)}]% ---------- Etapa 3 (v2): riel -16 V / 1,5 A (imagen especular de la etapa 2) ----------
% limitador con foldback: R_SC2, Darlington Q5 (TIP41C) + Q6 (BD139), sensor Q7 (BC547), R11-R14
% ampliación de corriente: R15 + Q8 (TIP41C); regulador U2 LM337 (Figs. 3-21, 3-26 y 3-27 de la monografía)
\draw (0,0) node[ocirc]{} node[left, net]{$V_{RAW-}$} -- (1.0,0) node[circ](a){} to[R, l={\lab{$R_{SC2}$}{\SI{0.27}{\ohm} 5 W}}] (2.8,0) node[circ](e1){} node[above right, font=\tiny]{$E_2$} -- (4.2,0);
\draw (4.2,0) node[npn, rotate=-90, xscale=-1, anchor=E](Q1){};
\node[lbl, anchor=south] at ($(Q1.B)+(0,1.05)$) {$Q_5$ TIP41C};
\draw (Q1.C) -- (7.0,0) node[circ](va){} node[above, font=\tiny]{$V_{A2}$};
% driver Q2 (BD140)
\coordinate (b1) at ($(Q1.B)+(0,-0.7)$);
\draw (Q1.B) -- (b1) node[circ]{};
\draw ($(b1)+(0,-0.5)$) node[npn, yscale=-1, anchor=E](Q2){};
\draw (b1) -- (Q2.E);
\node[lbl, anchor=west] at ($(Q2.C)+(0.15,0.35)$) {\lab{$Q_6$}{BD139}};
\draw (Q2.C) -- ++(0,-0.3) -| (7.0,-2.0); \draw (7.0,-2.0) node[circ]{};
\draw (3.6,0) node[circ]{} to[R, l={\lab{$R_{12}$}{\SI{1}{\kilo\ohm}}}] (3.6,-1.4) |- (b1);
\coordinate (bd) at ($(Q2.B)+(-0.4,0)$);
\draw (Q2.B) -- (bd) node[circ]{} node[above left, font=\tiny]{$BD_2$};
\draw (bd) -- (bd |- 0,-3.9) node[circ](bd2){};
\draw (bd2) to[R, l={\lab{$R_{11}$}{\SI{2.7}{\kilo\ohm} 1 W}}] (bd2 |- 0,-5.0) -- (bd2 |- 0,-6.2);
% sensor Q3 (BC557) y divisor de foldback R5-R6
\draw (1.0,-2.5) node[npn, xscale=-1, yscale=-1, anchor=E](Q3){};
\draw (a) -- (Q3.E);
\node[lbl, anchor=east] at ($(Q3.B)+(-0.95,0)$) {\lab{$Q_7$}{BC547}};
\draw (Q3.C) |- (bd2);
\draw (Q3.B) -- ++(0.4,0) node[circ](bs){} node[above right, font=\tiny]{$BS_2$};
\draw (bs) -- (bs |- 0,-1.0) to[R, l_={\lab{$R_{13}$}{\SI{270}{\ohm}}}] (2.8,-1.0) -- (e1);
\draw (bs) -- (bs |- 0,-5.4) -- (4.4,-5.4) to[R, l_={\lab{$R_{14}$}{\SI{12}{\kilo\ohm}}}] (6.6,-5.4) -- (7.0,-5.4) -- (va);
% ampliación de corriente
\draw (va) to[R, l_={\lab{$R_{15}$}{\SI{2.7}{\ohm} 1 W}}] (8.8,0) node[circ](in){} node[below right, font=\tiny]{$IN_2$} -- (10.2,0);
\draw (8.8,1.2) node[npn, rotate=-90, xscale=-1, anchor=B](Q4){};
\draw (in) -- (Q4.B);
\draw (Q4.E) -| (va);
\node[lbl, anchor=south] at ($(Q4.B)+(0,1.05)$) {$Q_8$ TIP41C};
\draw (9.6,0) node[circ]{} to[C, l_={\lab{$C_{10}$}{\SI{100}{\nano\farad}}}] (9.6,-6.2);
\reg{(10.2,0)}{U2 LM337T}{2}{3}{1}
\draw (12.8,0) -- (14.0,0) node[circ](out){};
\draw (Q4.C) -| (out);
\draw (out) |- (14.0,3.4) to[D*, invert, l_=$D_8$ 1N4007] (10.0,3.4) -- (10.0,0) node[circ]{};
\draw (11.5,-0.7) -- (11.5,-1.8) node[circ](adj){} node[above left, font=\tiny]{ADJ2};
\draw (adj) to[R, l_={\lab{$R_{17}$}{\SI{1.3}{\kilo\ohm} 1 \%}}] (11.5,-3.6) to[vR, l_={\lab{$RV_2$}{\SI{200}{\ohm} 3296}}] (11.5,-6.2);
\draw (adj) -- (12.9,-1.8) node[circ]{} to[eC, invert, l={\lab{$C_{11}$}{\SI{10}{\micro\farad}/25 V}}] (12.9,-6.2);
\draw (14.6,0) node[circ]{} to[R, l={\lab{$R_{16}$}{\SI{120}{\ohm} 1 \%}}] (14.6,-1.8) -- (12.9,-1.8);
\draw (14.6,-1.8) node[circ]{} -- (16.4,-1.8) to[D*, invert, l_={\lab{$D_9$}{1N4148}}] (16.4,0) node[circ]{};
\draw (out) -- (27.2,0) node[ocirc]{} node[right, net]{$-$16 V};
\draw (18.0,0) node[circ]{} to[eC, invert, l={\lab{$C_{12}$}{\SI{10}{\micro\farad} tánt.}}] (18.0,-6.2);
\draw (20.0,0) node[circ]{} to[C, l={\lab{$C_{13}$}{\SI{100}{\nano\farad}}}] (20.0,-6.2);
\draw (21.6,0) node[circ]{} to[eC, invert, l={\lab{$C_{14}$}{\SI{220}{\micro\farad}/25 V}}] (21.6,-6.2);
\draw (23.8,-6.2) to[D*, invert, l_={\lab{$D_{10}$}{1N5408}}] (23.8,0) node[circ]{};
\draw (25.2,0) node[circ]{} to[R, l={\lab{$R_{18}$}{\SI{2.2}{\kilo\ohm}}}] (25.2,-2.8) to[leDo, invert, l={\lab{LED2}{verde}}] (25.2,-6.2);
\draw (26.6,0) node[circ]{} -- (26.6,1.6) node[ocirc]{} node[left, font=\tiny]{J5-1};
\draw (26.6,-6.2) node[circ]{} -- (26.6,-7.2) node[ocirc]{} node[left, font=\tiny]{J5-2};
\draw[dashed] (26.6,1.6) -- (28.2,1.6) to[R, l={\lab{$R_L$ = 10 $\Omega$}{$\geq$30 W, externa}}] (28.2,-7.2) -- (26.6,-7.2);
\draw (0.4,-6.2) -- (27.2,-6.2) node[ocirc]{} node[right, net]{GND};
\foreach \x in {9.6,11.5,12.9,18.0,20.0,21.6,23.8,25.2} \draw (\x,-6.2) node[circ]{};
\draw (bd2 |- 0,-6.2) node[circ]{};
\draw (0.4,-6.2) node[ground]{};
\end{scope}
\node[font=\bfseries, anchor=west] at (-6.4,-18.0) {Riel +5 V (desde +16 V por $R_9$, regulador U3)};
\begin{scope}[shift={(-6.0,-20.4)}]% ---------- Etapa 4: regulador LM317 de +5 V alimentado desde +16 V, y bornera de salida ----------
\draw (0,0) node[ocirc]{} node[left, net]{+16 V} -- (0.6,0) to[R, l={\lab{$R_9$}{\SI{100}{\ohm} 3 W}}] (2.4,0) node[circ](vin){} node[below right, font=\tiny]{$V_{IN3}$} -- (7.0,0);
\draw (3.0,0) node[circ]{} to[eC, l={\lab{$C_{13}$}{\SI{100}{\micro\farad}/25 V}}] (3.0,-4);
\draw (5.5,0) node[circ]{} to[C, l={\lab{$C_{14}$}{\SI{100}{\nano\farad}}}] (5.5,-4);
\reg{(7.0,0)}{U3 LM317T}{3}{2}{1}
\draw (9.6,0) -- (10.4,0) node[circ](out){};
\draw (out) -- (10.4,1.4) to[D*, l_=$D_9$ 1N4007] (6.6,1.4) -- (6.6,0) node[circ]{};
\draw (8.3,-0.7) -- (8.3,-1.6) node[circ](adj){} node[above left, font=\tiny]{ADJ3};
\draw (adj) to[R, l_={\lab{$R_{11}$}{\SI{330}{\ohm} 1 \%}}] (8.3,-2.8) to[vR, l_={\lab{$RV_3$}{\SI{50}{\ohm} 3296}}] (8.3,-4);
\draw (adj) -- (9.6,-1.6) node[circ]{} to[eC, l={\lab{$C_{15}$}{\SI{10}{\micro\farad}/25 V}}] (9.6,-4);
\draw (out) -- (11.2,0) node[circ]{} to[R, l={\lab{$R_{10}$}{\SI{120}{\ohm} 1 \%}}] (11.2,-1.6) -- (9.6,-1.6);
\draw (11.2,-1.6) node[circ]{} -- (13.2,-1.6) to[D*, l_={\lab{$D_{10}$}{1N4148}}] (13.2,0) node[circ]{};
\draw (out) -- (21.0,0) node[ocirc]{} node[right, net]{+5 V};
\draw (15.0,0) node[circ]{} to[eC, l={\lab{$C_{16}$}{\SI{10}{\micro\farad} tánt.}}] (15.0,-4);
\draw (16.7,0) node[circ]{} to[C, l={\lab{$C_{17}$}{\SI{100}{\nano\farad}}}] (16.7,-4);
\draw (18.3,-4) to[D*, l_={\lab{$D_{11}$}{1N4007}}] (18.3,0) node[circ]{};
\draw (19.9,0) node[circ]{} to[R, l={\lab{$R_{14}$}{\SI{470}{\ohm}}}] (19.9,-2.2) to[leDo, l={\lab{LED3}{verde}}] (19.9,-4);
\draw (0.6,-4) -- (21.0,-4) node[ocirc]{} node[right, net]{GND};
\foreach \x in {3.0,5.5,8.3,9.6,15.0,16.7,18.3,19.9} \draw (\x,-4) node[circ]{};
\draw (0.6,-4) node[ground]{};
\end{scope}
\node[font=\bfseries, anchor=west] at (-6.4,-25.6) {Bornera de salida y puntos de prueba};
\begin{scope}[shift={(-3.0,-26.6)}]% ---------- Bornera de salida J2 y puntos de prueba ----------
\draw[thick] (0,0) rectangle (2.2,-4.2); \node[font=\footnotesize\bfseries] at (1.1,0.35) {J2 (bornera 4 vías, 5,08 mm)};
\foreach \y/\n/\t in {-0.6/1/+16 V, -1.6/2/$-$16 V, -2.6/3/GND, -3.6/4/+5 V} {
  \draw (-1.6,\y) node[left, net]{\t} -- (0,\y); \node[font=\scriptsize] at (0.4,\y) {\n};
  \draw (2.2,\y) -- (3.6,\y) node[ocirc]{};
}
\node[font=\scriptsize, align=left, anchor=west] at (3.8,-2.1) {hacia J\_PWR de las tarjetas PID\\(pin 1 +16 V, 2 $-$16 V,\\3 GND, 4 +5 V: mismo orden)};
\draw[thick] (8.6,0) rectangle (10.6,-3.2); \node[font=\footnotesize\bfseries] at (9.6,0.35) {J3 (regleta 1$\times$3)};
\foreach \y/\n/\t in {-0.6/1/$V_{RAW+}$ (TP1), -1.6/2/GND, -2.6/3/$V_{RAW-}$ (TP2)} {
  \draw (10.6,\y) -- (11.2,\y) node[right, font=\scriptsize]{\t}; \node[font=\scriptsize] at (8.9,\y) {\n};
}
\node[font=\scriptsize, align=left, anchor=north west] at (8.6,-3.4) {puntas del osciloscopio para medir\\el rizado de $C_1$ y $C_{2A}$/$C_{2B}$};
\node[font=\scriptsize, align=left, anchor=north west] at (0,-4.8) {J4 (+16 V / GND) y J5 ($-$16 V / GND): borneras de 2 vías para conectar\\y quitar la resistencia de carga de prueba $R_L$ = 10 $\Omega$ (ver etapas 2 y 3).};
\end{scope}
\end{circuitikz}
\end{document}
